% Tres perfiles de 2025-10-30.pdf P01: raíz simple, doble y triple en cero.
% Ejemplos ilustrativos; no son gráficas de polinomios especificados allí.

\begin{tikzpicture}[x=1.35cm,y=1.15cm,line width=.6pt,font=\small]
\foreach \x/\y/\m in {0/0/1,5/0/2,2.5/-3.5/3}{
  \begin{scope}[shift={(\x,\y)}]
  \draw[grisCIMAT] (-.9,0)--(2.05,0);
  \draw[grisCIMAT] (0,-1)--(0,1.35);
  \node[below] at (.65,-1.15) {Multiplicidad $\m$};
  \end{scope}
}
\begin{scope}
\clip (-.9,-1) rectangle (2.05,1.35);
\draw[black,samples=100,domain=-.9:2.05,smooth] plot (\x,{.8*\x*(\x-1.6)});
\end{scope}
\begin{scope}[shift={(5,0)}]
\clip (-.9,-1) rectangle (2.05,1.35);
\draw[black,samples=100,domain=-.9:2.05,smooth] plot (\x,{3*\x*\x*(1.2-\x)});
\end{scope}
\begin{scope}[shift={(2.5,-3.5)}]
\clip (-.9,-1) rectangle (2.05,1.35);
\draw[black,samples=100,domain=-.9:2.05,smooth] plot (\x,{\x*\x*\x*(\x-1.5)});
\end{scope}
\end{tikzpicture}
